\documentclass[letterpaper,10pt]{article}
\usepackage[margin=1cm]{geometry}
\usepackage{graphicx}
\usepackage{xcolor}
\usepackage{tikz}
\usepackage{circuitikz}
\usetikzlibrary{arrows.meta,patterns,calc}
\definecolor{azulFuerte}{HTML}{1A5276}
\pagestyle{empty}

\newcommand{\comparar}[3]{%
    \section*{#1}
    \begin{minipage}[t]{0.48\textwidth}
        \centering
        \textbf{Original PNG}\\[0.2cm]
        \includegraphics[width=\textwidth,height=0.34\textheight,keepaspectratio]{#2}
    \end{minipage}\hfill
    \begin{minipage}[t]{0.48\textwidth}
        \centering
        \textbf{Nuevo CircuitikZ}\\[0.2cm]
        \resizebox{\textwidth}{!}{\input{#3}}
    \end{minipage}
    \newpage
}

\begin{document}
\comparar{Ejercicio 2}{Imagenes Ejercicios/Ejercicio_2.png}{circuits/ejercicio_2.tex}
\comparar{Ejercicio 3}{Imagenes Ejercicios/Ejercicio_3.png}{circuits/ejercicio_3.tex}
\comparar{Ejercicio 4}{Imagenes Ejercicios/Ejercicio_4.png}{circuits/ejercicio_4.tex}
\comparar{Ejercicio 5.1}{Imagenes Ejercicios/Ejercicio_5_1.png}{circuits/ejercicio_5_1.tex}
\comparar{Ejercicio 5.2}{Imagenes Ejercicios/Ejercicio_5_2.png}{circuits/ejercicio_5_2.tex}
\comparar{Ejercicio 6}{Imagenes Ejercicios/Ejercicio_6.png}{circuits/ejercicio_6.tex}
\end{document}
